%%%%%%%%%%%%%%%%%%%%%%%%%%%%%%%%%%%%%%%%%%%%%%%%%%%%%%%%%%%%
%% Lecture 12
%%%%%%%%%%%%%%%%%%%%%%%%%%%%%%%%%%%%%%%%%%%%%%%%%%%%%%%%%%%%

\lecture
{12}
{Wednesday, 18 September 2026}
{Introduction to Probability}
{Dr. Nishant Chandgotia}


%============================================================
% Convolution of Measures
%============================================================

\begin{observation}

Suppose that
\(
\mu=f\,dm,
\qquad
\nu=g\,dm.
\)

Then the product measure is
\(
\mu\times\nu,
\)
and the convolution of $\mu$ and $\nu$ has density
\(
f*g.
\)

More precisely,
\[
\mu*\nu=(f*g)\,dm,
\]
where
\[
(f*g)(z)
=
\int_{\mathbb{R}}
f(z-y)g(y)\,dy.
\]

A proof of this result can be found in Durrett.

\end{observation}


%============================================================
% Convolution with Counting Measure
%============================================================

\begin{observation}
If the underlying measure is the counting measure, then the
convolution takes the simple form
\[
(f*g)(n)
=
\sum_{m}f(m)g(n-m).
\]

Thus, in the discrete case, convolution is given by a sum rather
than an integral.

\end{observation}

\section{ Distribution Function of a Sum
}


%============================================================
% Distribution Function of a Sum
%============================================================

\begin{theorem}{Distribution Function of a Sum}
{cdf-sum-product-measure}

Let
\(
X\sim\mu,
\qquad
Y\sim\nu,
\)
where $X$ and $Y$ are independent.

If $F$ and $G$ are the distribution functions of $\mu$ and $\nu$,
respectively, then the distribution function of $X+Y$ is
\[
\begin{aligned}
\mathbb{P}(X+Y\leq z)
&=
\mathbb{P}\bigl((X,Y)\in
\{(x,y):x+y\leq z\}\bigr)\\
&=
\int_{\{x+y\leq z\}}
d(\mu\times\nu)(x,y).
\end{aligned}
\]

Equivalently,
\[
\mathbb{P}(X+Y\leq z)
=
\int_{\mathbb{R}}F(z-y)\,dG(y).
\]

\end{theorem}

\section{Some Questions and Exercises on Sum and Product of Random Variables}

%============================================================
% Exercises
%============================================================

\begin{question}{Density of the Product of Two Independent Random Variables}
{density-product-random-variables}

Let
\(
X\sim\mu,
\qquad
Y\sim\nu,
\)
where
\(
\mu=f\,dm,
\qquad
\nu=g\,dm.
\)

Assume that $X$ and $Y$ are independent.
Compute the density of
\(
XY.
\)

\end{question}


\begin{question}{Sum of Two Independent Poisson Random Variables}
{sum-independent-poisson}

Let
\(
X\sim\operatorname{Poi}(\lambda),
\qquad
Y\sim\operatorname{Poi}(\mu),
\)
and suppose that $X$ and $Y$ are independent.

Determine the distribution of
\(
X+Y.
\)

\end{question}


\begin{question}{Sum of Two Independent Binomial Random Variables}
{sum-independent-binomial}

Let
\(
X\sim\operatorname{Bin}(m,p),
\qquad
Y\sim\operatorname{Bin}(m,p),
\)
and suppose that $X$ and $Y$ are independent.

Determine the distribution of
\(
X+Y.
\)

\end{question}


%============================================================
% Sum of Independent Exponential Random Variables
%============================================================

\begin{example}{Sum of Two Independent Exponential Random Variables}
{sum-two-exponential}

Suppose
\(
X,Y\sim\operatorname{Exp}(\lambda)
\)
are independent, where $\lambda>0$ is the rate parameter.

The densities of $X$ and $Y$ are
\[
f_X(x)=\lambda e^{-\lambda x},
\qquad x>0,
\]
and
\[
f_Y(y)=\lambda e^{-\lambda y},
\qquad y>0.
\]

Let
\(
Z=X+Y.
\)

By the convolution formula,
\[
f_Z(y)
=
\int_0^y
f_X(x)f_Y(y-x)\,dx.
\]

Therefore,
\[
\begin{aligned}
f_Z(y)
&=
\int_0^y
\lambda e^{-\lambda x}
\lambda e^{-\lambda(y-x)}
\,dx\\
&=
\lambda^2e^{-\lambda y}
\int_0^y dx\\
&=
\lambda^2y e^{-\lambda y},
\qquad y>0.
\end{aligned}
\]

Thus
\(
X+Y
\)
has a Gamma distribution with shape parameter $2$ and rate
parameter $\lambda$:
\(
X+Y\sim\Gamma(2,\lambda).
\)

\end{example}


%============================================================
% Sum of n Independent Exponential Random Variables
%============================================================

\begin{proposition}{Sum of Independent Exponential Random Variables}
{sum-independent-exponential}

Let
\(
X_1,X_2,\ldots,X_n
\)
be independent random variables such that
\(
X_i\sim\operatorname{Exp}(\lambda),
\qquad
\lambda>0.
\)

Then
\[
X_1+\cdots+X_n
\sim
\Gamma(n,\lambda).
\]

Its density is
\[
f(x)
=
\frac{\lambda^n}{(n-1)!}
x^{n-1}e^{-\lambda x},
\qquad x>0.
\]

\end{proposition}

\begin{proof}

We prove the result by induction on $n$.

For $n=1$,
\[
f_{X_1}(x)
=
\lambda e^{-\lambda x},
\qquad x>0.
\]

Since
\(
Gamma(1,\lambda)
\)
has density
\(
\frac{\lambda^1}{\Gamma(1)}
x^{1-1}e^{-\lambda x}
=
\lambda e^{-\lambda x},
\)
the result holds for $n=1$.

\medskip

Suppose the result holds for $n$. Then
\(
S_n=X_1+\cdots+X_n
\)
has density
\[
f_{S_n}(x)
=
\frac{\lambda^n}{(n-1)!}
x^{n-1}e^{-\lambda x},
\qquad x>0.
\]

Since $X_{n+1}$ is independent of $S_n$,
\(
S_{n+1}=S_n+X_{n+1}
\)
has density given by convolution:
\[
\begin{aligned}
f_{S_{n+1}}(x)
&=
\int_0^x
f_{S_n}(u)
f_{X_{n+1}}(x-u)\,du\\
&=
\int_0^x
\frac{\lambda^n}{(n-1)!}
u^{n-1}e^{-\lambda u}
\lambda e^{-\lambda(x-u)}
\,du\\
&=
\frac{\lambda^{n+1}e^{-\lambda x}}
{(n-1)!}
\int_0^x u^{n-1}\,du\\
&=
\frac{\lambda^{n+1}e^{-\lambda x}}
{(n-1)!}
\frac{x^n}{n}\\
&=
\frac{\lambda^{n+1}}
{n!}
x^ne^{-\lambda x}.
\end{aligned}
\]

Thus
\(
S_{n+1}
\)
has density
\(
f_{S_{n+1}}(x)
=
\frac{\lambda^{n+1}}{n!}
x^ne^{-\lambda x},
\qquad x>0.
\)

Therefore, by induction,
\[
\boxed{
X_1+\cdots+X_n
\sim
\Gamma(n,\lambda).
}
\]

\end{proof}


%============================================================
% Gamma Distribution
%============================================================

\begin{definition}{Gamma Distribution}{gamma-distribution}

Let
\(
\alpha>0,
\qquad
\lambda>0.
\)


A random variable $X$ is said to have a Gamma distribution with
shape parameter $\alpha$ and rate parameter $\lambda$ if its
density is
\[
f_X(x)
=
\frac{\lambda^\alpha}{\Gamma(\alpha)}
x^{\alpha-1}e^{-\lambda x},
\qquad x>0,
\]
and
\[
f_X(x)=0,
\qquad x\leq0.
\]

We write
\[
X\sim Gamma(\alpha,\lambda).
\]

The Gamma function is defined by
\[
\Gamma(\alpha)
=
\int_0^\infty
x^{\alpha-1}e^{-x}\,dx,
\qquad
\alpha>0.
\]

For $n\in\mathbb{N}$,
\[
\Gamma(n)=(n-1)!.
\]

\end{definition}


%============================================================
% Sum of Independent Gamma Random Variables
%============================================================

\begin{proposition}{Sum of Independent Gamma Random Variables}
{sum-independent-gamma}

Let
\[
X\sim Gamma(\alpha,\lambda),
\qquad
Y\sim Gamma(\beta,\lambda),
\]
where $X$ and $Y$ are independent and $\lambda$ is the common
rate parameter.

Then
\[
\boxed{
X+Y\sim Gamma(\alpha+\beta,\lambda).
}
\]

\end{proposition}

\begin{proof}

Let
\[
Z=X+Y.
\]

Since $X$ and $Y$ are independent, the density of $Z$ is given by
the convolution of the densities of $X$ and $Y$:
\[
f_Z(z)
=
\int_{-\infty}^{\infty}
f_X(x)f_Y(z-x)\,dx.
\]

Since $X>0$ and $Y>0$, for $z>0$ we must have
\(
x>0
\qquad\text{and}\qquad
z-x>0.
\)

Thus
\[
0<x<z.
\]

Therefore,
\[
f_Z(z)
=
\int_0^z
f_X(x)f_Y(z-x)\,dx.
\]

Substituting the Gamma densities,
\[
\begin{aligned}
f_Z(z)
&=
\int_0^z
\frac{\lambda^\alpha}{\Gamma(\alpha)}
x^{\alpha-1}e^{-\lambda x}
\frac{\lambda^\beta}{\Gamma(\beta)}
(z-x)^{\beta-1}e^{-\lambda(z-x)}
\,dx\\
&=
\frac{\lambda^{\alpha+\beta}}
{\Gamma(\alpha)\Gamma(\beta)}
\int_0^z
x^{\alpha-1}(z-x)^{\beta-1}
e^{-\lambda x}e^{-\lambda(z-x)}
\,dx.
\end{aligned}
\]

But
\[
e^{-\lambda x}e^{-\lambda(z-x)}
=
e^{-\lambda z}.
\]

Hence
\[
f_Z(z)
=
\frac{\lambda^{\alpha+\beta}e^{-\lambda z}}
{\Gamma(\alpha)\Gamma(\beta)}
\int_0^z
x^{\alpha-1}(z-x)^{\beta-1}\,dx.
\]

We now evaluate the integral
\[
I
=
\int_0^z
x^{\alpha-1}(z-x)^{\beta-1}\,dx.
\]

Make the substitution
\(
u=\frac{x}{z}.
\)

Then
\(
x=zu,
\qquad
dx=z\,du,
\)
and
\(
z-x=z(1-u).
\)

When $x=0$, $u=0$, and when $x=z$, $u=1$. Therefore,
\[
\begin{aligned}
I
&=
\int_0^1
(zu)^{\alpha-1}
\bigl(z(1-u)\bigr)^{\beta-1}
z\,du\\
&=
z^{\alpha+\beta-1}
\int_0^1
u^{\alpha-1}(1-u)^{\beta-1}\,du.
\end{aligned}
\]

Recall the Beta function:
\[
B(\alpha,\beta)
=
\int_0^1
u^{\alpha-1}(1-u)^{\beta-1}\,du.
\]

The Beta--Gamma identity gives
\[
B(\alpha,\beta)
=
\frac{\Gamma(\alpha)\Gamma(\beta)}
{\Gamma(\alpha+\beta)}.
\]

Consequently,
\[
I
=
z^{\alpha+\beta-1}
\frac{\Gamma(\alpha)\Gamma(\beta)}
{\Gamma(\alpha+\beta)}.
\]

Substituting this into the expression for $f_Z(z)$,
\[
\begin{aligned}
f_Z(z)
&=
\frac{\lambda^{\alpha+\beta}e^{-\lambda z}}
{\Gamma(\alpha)\Gamma(\beta)}
z^{\alpha+\beta-1}
\frac{\Gamma(\alpha)\Gamma(\beta)}
{\Gamma(\alpha+\beta)}\\
&=
\frac{\lambda^{\alpha+\beta}}
{\Gamma(\alpha+\beta)}
z^{\alpha+\beta-1}e^{-\lambda z},
\qquad z>0.
\end{aligned}
\]

Also,
\[
f_Z(z)=0,
\qquad z\leq0.
\]

But this is precisely the density of
\[
Gamma(\alpha+\beta,\lambda).
\]

Therefore,
\[
\boxed{
X+Y\sim Gamma(\alpha+\beta,\lambda).
}
\]

\end{proof}


%============================================================
% Final Note
%============================================================



Recall the Carathéodory Extension Theorem.



%============================================================
% End of Lecture
%============================================================

\lectureend
